\documentclass[11pt,a4paper]{article}
\usepackage{fontspec}
\setmainfont{Arial}
\setsansfont{Arial}
\usepackage{polyglossia}
\setdefaultlanguage{english}
\usepackage{geometry,xcolor,hyperref,fancyhdr,enumitem,tabularx,booktabs}
\geometry{margin=18mm}
\setlength{\parindent}{0pt}
\setlength{\parskip}{5pt}
\setlist{nosep,leftmargin=*}
\definecolor{navy}{HTML}{17365D}\definecolor{teal}{HTML}{167D8D}\definecolor{pale}{HTML}{EEF5F7}\definecolor{amber}{HTML}{B45309}
\setlength{\headheight}{14pt}
\pagestyle{fancy}\fancyhf{}\lhead{MATEMAFIJA}\rhead{Final Five-Minute Pitch}\cfoot{\thepage}
\newcommand{\speaker}[2]{\vspace{1mm}\textcolor{teal}{\textbf{#1 $\mid$ #2}}\par}
\newcommand{\cue}[1]{\textcolor{amber}{\textit{[#1]}}}
\newcommand{\boxtext}[1]{\vspace{2mm}\noindent\fcolorbox{teal}{pale}{\parbox{0.94\linewidth}{#1}}\vspace{2mm}}
\title{\vspace{-1cm}\textbf{MATEMAFIJA}\\[2mm]\Large Final Five-Minute Speaking Script\\[2mm]\large 3 minutes story + 2 minutes prototype demo}
\author{Galina and Gleb --- VILNIUS TECH}
\date{Hack4Vilnius 2026 --- final repository state}

\begin{document}
\maketitle
\boxtext{Target: calm delivery, approximately 650 spoken words including demo narration. Cues in brackets are not spoken. Galina owns the story and city/business value; Gleb owns the live evidence and mechanism demo. MCU, contributions and response effects are simulated unless explicitly labelled observed.}

\section*{Slide 1 --- Team and hook}
\speaker{GALINA}{0:00--0:18}
Imagine it is a weekday evening in Vilnius. Your car needs charging before tomorrow, but your usual station has no available charging point. Do you wait, keep searching, or drive somewhere else? We are Galina and Gleb, Applied Mathematics students from VILNIUS TECH, and this is MATEMAFIJA.

\section*{Slide 2 --- The city problem}
\speaker{GALINA}{0:18--0:52}
Vilnius already has charging infrastructure and public data, but availability is uneven. Our current collection contains more than 2.8 million OCPI status observations and over 2,100 current EVSE records at locations marked Vilnius. Yet the data do not show a physical queue, individual waiting time or why a driver leaves. So simply inventing a congestion fee would be easy --- and scientifically weak.

\section*{Slide 3 --- The solution}
\speaker{GALINA}{0:52--1:35}
We built a policy stress-testing digital twin. First, it reconstructs usable charging capacity from OCPI and DATEX evidence. OSA-v1 activates a small, capped scarcity contribution only when low availability is qualified and sustained. Prior funded contributions can then support Mobility Credit Unit offers at suitable, less-scarce alternatives. The user sees a positive option, not only a restriction. The model asks whether this complete contribution-to-credit flow improves availability without suppressing accepted use, overloading another station or spending the same credit twice.

\section*{Slide 4 --- Why it is credible and useful}
\speaker{GALINA}{1:35--2:25}
The key is that we separate facts from assumptions. Statuses, timestamps and tariffs are observed. Availability intervals and scarcity are reconstructed. Price, credit, distance and delay sensitivities are explicitly assumed because this policy has not been deployed. Fixed and adaptive scenarios use the same baseline and random draws. Every result must pass four gates: system improvement, safe redistribution, an observable user-distance proxy and exact accounting reconciliation. If evidence is insufficient, the intervention is withheld --- not silently treated as zero.

The daily beneficiaries are drivers. Charging operators and the city are the first potential customers of the policy-testing and pilot-design tool. Our headline pilot KPI is simple: reduce qualified zero-availability without increasing scarcity at receiver stations.

\section*{Slide 5 --- Demo: observed evidence}
\speaker{GLEB}{2:25--3:18}
\cue{Open Evidence. Show the map, source timestamp and readiness panel.}
Here is the live evidence layer. These are source records, not a mock dataset. We preserve the OCPI hierarchy, tariff components and timezone-aware history. For each site we separate available, charging and equipment-related states, because broken infrastructure must not trigger a demand policy. The collection currently spans just over four days. Our formal transition gate requires seven, so the effectiveness heatmap is deliberately locked. That lock is a feature: the prototype refuses to present assumptions as evidence.

\section*{Slide 6 --- Demo: policy, response and credit flow}
\speaker{GLEB}{3:18--4:25}
\cue{Open Policy, then Response. Move one credit or distance control only.}
When evidence qualifies, OSA-v1 produces a bounded signal after a fifteen-minute grace period. It is not an actual bill. The Policy view connects that signal to a limited MCU pool and funded receiver offers. In Response, we vary release sensitivity, accepted-start sensitivity, credit value, distance and delay. These curves are hypothetical mechanisms, not measured Vilnius elasticity.

\cue{Show the flow and four evaluation cards.}
A selected alternative reserves one capacity unit and the required MCU atomically. Only a simulated receiver start may issue the credit. Offer, reservation, issuance and redemption remain separate. We then compare the fixed and adaptive arms across system, redistribution, user-proxy and accounting outputs. No single good scarcity number can hide displaced demand or a failed ledger.

\section*{Slide 7 --- Benefit, business path and ask}
\speaker{GALINA}{4:25--5:00}
MATEMAFIJA gives Vilnius something more useful than an untested pricing promise: it shows how strong the response must be, where the trade-offs appear and what a pilot must measure. Our next step is a bounded city--operator pilot with selected sites, Session or CDR evidence and pseudonymous choice records. That would replace assumed sensitivities with measured Vilnius behaviour.

Lithuania's DAEI system is a relevant accounting precedent and a possible future operator-level bridge, but MCU and DAEI are separate today. We are not selling a finished payment system. We are offering a transparent way to decide whether one should be piloted. MATEMAFIJA --- turning scarcity into better charging choices. Thank you.

\newpage
\section*{Seven-slide production plan}
\begin{tabularx}{\textwidth}{p{1.0cm}p{4.0cm}X}
\toprule & Slide & Visual proof \\\midrule
1 & Team + hook & Vilnius evening, ``No available point --- what now?'' \\
2 & Problem & Marta/origin situation plus one live-data number \\
3 & Solution & Observe $\to$ qualify scarcity $\to$ fund MCU $\to$ offer alternative \\
4 & Credibility & observed/reconstructed/assumed/simulated and four gates \\
5 & Evidence demo & real map, source timestamp, tariffs and readiness lock \\
6 & Mechanism demo & OSA signal, response control, MCU/capacity flow and evaluation cards \\
7 & Benefit + ask & Drivers / CPO / Vilnius, one KPI, bounded pilot, DAEI future path \\
\bottomrule
\end{tabularx}

\section*{Recording rules}
\begin{itemize}
\item Rehearse the full clicks and speech at least three times; target 4:45--4:55.
\item Keep only one live control movement. Preload the database; never wait for network collection.
\item Record a clean demo clip and save screenshots of Evidence, Policy, Response and Decision.
\item Update the rounded observation count only if it remains easy to say and verify.
\item Never call zero availability a queue; never call MCU money, cashback or DAEI.
\item Submit the video as Public or Unlisted and verify the link in an incognito window.
\end{itemize}

\section*{Fast jury answers}
\textbf{What is innovative?} We test a complete, budget-constrained scarcity-to-credit circulation mechanism, including behavioural uncertainty, receiver spillover and accounting --- not only a fee formula.

\textbf{Who uses and pays?} Drivers benefit; a CPO and the city use the policy lab. The first commercial hypothesis is a city--operator pilot and decision-support engagement. Real funding and liability must be contracted before credits carry value.

\textbf{Can you prove behaviour changes?} No. We expose assumed response parameters and define the minimum response that a pilot must validate.

\textbf{Why is the heatmap locked?} Four days of data do not meet the seven-day transition gate. Showing effectiveness now would disguise assumptions as evidence.

\textbf{Are credits real?} No. MCU is an integer simulation unit with reconciled lifecycle states. The prototype processes no payment and holds no personal wallet.

\textbf{Is this DAEI?} No. DAEI is a regulated renewable-transport accounting system. It is a precedent and possible future integration path, not the present MCU mechanism.

\textbf{What do you need next?} Qualified pilot sites, an operator partner, Session/CDR evidence and privacy-safe choice/redemption measurements.

\vfill
\small\textit{Final verbal claim: Under declared response assumptions, MATEMAFIJA identifies when a contribution-to-credit policy is jointly robust within observable system, redistribution, user-proxy and accounting limits, and turns the uncertainty into a measurable Vilnius pilot.}
\end{document}
